% C1 FlashKDA 全轮次实验记录
% 编译：latexmk -xelatex -interaction=nonstopmode C1_EXPERIMENT_RECORD.tex
% 这是过程记录，不是答辩稿。集群目录中的 <user>/<repo> 是脱敏占位符。
\documentclass[11pt,a4paper,fontset=none]{ctexart}
\setCJKmainfont[BoldFont={Noto Serif CJK SC Bold}]{Noto Serif CJK SC}
\setCJKsansfont[BoldFont={Noto Sans CJK SC Bold}]{Noto Sans CJK SC}
\setCJKmonofont{Noto Sans CJK SC}
\usepackage{geometry}
\geometry{margin=2.3cm}
\usepackage{amsmath,amssymb}
\usepackage{booktabs,array,multirow,longtable,tabularx}
\newcolumntype{L}[1]{>{\raggedright\arraybackslash}p{#1}}
\usepackage{enumitem}
\setlist{nosep,leftmargin=2em}
\usepackage{xcolor}
\definecolor{accent}{RGB}{28,60,120}
\definecolor{codebg}{RGB}{248,248,246}
\usepackage[most]{tcolorbox}
\usepackage{listings}
\definecolor{strgreen}{RGB}{21,101,42}
\lstset{
  basicstyle=\ttfamily\small,
  backgroundcolor=\color{codebg},
  frame=single, rulecolor=\color{gray!40}, framerule=0.4pt,
  breaklines=true, columns=fullflexible, keepspaces=true,
  showstringspaces=false, language=bash,
  morekeywords={make,nvcc,uv,pytest,python,cd,ncu,nvdisasm,cuobjdump,srun,sinfo,squeue,ssh,scp},
  commentstyle=\color{gray!130}\itshape,
  keywordstyle=\color{accent}\bfseries,
  stringstyle=\color{strgreen},
  xleftmargin=4pt, xrightmargin=4pt,
}
\usepackage{hyperref}
\hypersetup{colorlinks, linkcolor=accent, urlcolor=accent}
\Urlmuskip=0mu plus 1mu
\emergencystretch=3em
\pagestyle{plain}
\setlength{\LTpre}{0.6em}
\setlength{\LTpost}{0.6em}

\newtcolorbox{record}[1]{enhanced,breakable,
  colback=gray!5,colframe=gray!50,boxrule=0.5pt,arc=1.5pt,
  left=7pt,right=7pt,top=5pt,bottom=5pt,
  title={#1},fonttitle=\bfseries\sffamily\small,coltitle=black,
  attach title to upper={\quad}}
\newtcolorbox{decision}[1]{enhanced,breakable,
  colback=accent!4,colframe=accent!70,boxrule=0.7pt,arc=1.5pt,
  left=7pt,right=7pt,top=5pt,bottom=5pt,
  title={#1},fonttitle=\bfseries\sffamily}
\newcommand{\FlashKDA}{\textsc{FlashKDA}}
\newcommand{\Fla}{\textsc{FLA}}
\newcommand{\code}[1]{\texttt{#1}}

\begin{document}
\begin{center}
  {\LARGE\bfseries C1：FlashKDA 全轮次实验记录}\\[6pt]
  {\sffamily 作业 2 Team / 可审计实验日志（不是答辩版本）}
\end{center}
\vspace{0.5em}\hrule\vspace{1em}

\begin{record}{文档范围与使用说明}
本文是 C1 从第一轮起持续追加的实验底稿，而不是一次性的答辩稿。第一轮记录
了题目与源码阅读、两套 Slurm 集群的接入和环境构建、FlashKDA extension 的
编译、官方正确性和 benchmark、NCU/SASS 检查、SM100 tcgen05 正反例、CHUNK
数值范围 probe 与 state 精度 probe。第二轮及以后每轮均沿用相同的「动机
--假设--命令--原始观察--解释--局限」结构追加，不能用后来的结论改写早期
观察。文中的命令均来自实际执行过程；为避免把个人用户名、IP 和绝对目录写进
git，使用 \code{/home/lcpu/<user>/<repo>} 和环境变量占位符。

轻量日志已经版本控制在 \code{results/}。完整 NCU report 约 90\,MiB，
两份 SASS dump 也较大，按 \code{.gitignore} 仅保留在实验运行环境；
关键命令、计数和指标均在本文中展开记录。
\end{record}

\begin{record}{轮次与修订约定}
本文件的文件名固定为 \code{C1\_EXPERIMENT\_RECORD.tex}；每一次有意义的
实验进展在末尾新开「轮次」章节，并在该章内写明运行日期、代码 commit、设备、
环境、命令、输入和局限。若旧数据在后续被复核，应在新章节中说明复核，而不
静默修改旧章节。PDF 是本地编译产物，不随 git 提交；可审计的源文件、轻量日志
和脚本才进入版本控制。

\begin{center}
\begin{tabular}{@{}lll@{}}
\toprule
轮次 & 覆盖范围 & 当前状态 \\
\midrule
第一轮 & 环境、基线、指令/数值探针和 tcgen05 可行性 & 已完成并冻结为里程碑 \\
第二轮 & 5090 误差诊断、K2 并行度设计与 prototype & 进行中 \\
后续轮次 & 大 CHUNK/rescale 与 SM100 v2 决策收束 & 待开始 \\
\bottomrule
\end{tabular}
\end{center}
\end{record}

\tableofcontents
\clearpage

\section{题目、范围与轮次路线图}

\subsection{C1 要回答的问题}

C1 的题面是「FlashKDA---官方 kernel 停在 SM80 MMA」。目标不仅是让项目
运行，而是建立或推翻如下解释：在 B300/GB200 等 Blackwell 上，
\FlashKDA{} 为什么依然使用 SM80 风格 \code{mma.sync}/HMMA，而不使用
SM90 WGMMA 或 SM100 \code{tcgen05}。

题面给出三个层次：第一层在 B300 上复现并用 NCU/SASS 验证；第二层回答
CHUNK、tcgen05 tile、递推并行度、瓶颈、state 精度和 SM100 专版决策六个
讨论点；第三层从「只换指令」「大 CHUNK + rescale」「并行度重构」任选一条
动手。题面允许负结果，但必须用代码、数据和设计分析说明负在哪里。

\subsection{第一轮的策略}

第一轮刻意不修改生产 \FlashKDA{} kernel。先建立可复现的正确性、性能和
instruction 基线；之后选风险最低且最能排除错误直觉的挑战切面：
验证 tcgen05 是否可直接替换 \(16\times16\) SM80 MMA。并行度重构和大 CHUNK
都会同时改变数据布局、数值和性能，不适合在没有基线时盲改。

\begin{longtable}{@{}L{0.13\linewidth}L{0.32\linewidth}L{0.34\linewidth}L{0.13\linewidth}@{}}
\toprule
编号 & 子问题 & 第一轮方法 & 状态 \\
\midrule
\endhead
\bottomrule
\endlastfoot
P0 & 两台集群能否构建与运行？ & SSH/Slurm probe；GPU allocation 内检查 & 完成 \\
P1 & 生产路径是否仍是 SM80？ & 源码定位、NCU、SASS 计数 & 完成 \\
P2 & B300 正确性是否通过？ & 官方 fixed、varlen、FLA 对拍 & 完成 \\
P3 & H=96 基线速度是多少？ & 官方 30/200/5 event benchmark & 完成 \\
P4 & K1/K2 的资源形态？ & NCU 指标和源码资源映射 & 初步完成 \\
P5 & tcgen05 能否直接换掉 16x16 atom？ & CUTLASS 正例 + 16x16 编译负例 & 完成 \\
P6 & 大 CHUNK 最先坏在哪里？ & gate cumulative exponent bf16 range probe & 初步完成 \\
P7 & fp32 state 是否内部真为 fp32？ & 源码路径 + 三方精度 probe & 初步完成 \\
P8 & 如何增加 K2 并行度？ & 候选方案和资源上界 & 第二轮 \\
\end{longtable}

\section{代码树、版本和源码阅读记录}

\subsection{目录和版本}

\begin{longtable}{@{}L{0.29\linewidth}L{0.24\linewidth}L{0.37\linewidth}@{}}
\toprule
路径 & 角色 & 第一轮使用方式 \\
\midrule
\endhead
\bottomrule
\endlastfoot
\code{TASK.md} & C1 题面 & 三层任务、官方 shape、讨论点和验收口径 \\
\code{FlashKDA/} & Moonshot 快照 & CUDA kernel、测试、benchmark、设计文档 \\
\code{fla\_kda\_ref/} & FLA reference 快照 & PyTorch/Triton 对照与接口语义 \\
\code{challenge/} & 第一轮新增 & tcgen05、range/state probe、构建脚本 \\
\code{results/} & 第一轮证据 & 正确性、benchmark、失败 log、probe 与 NCU CSV \\
\end{longtable}

\FlashKDA{} 快照对应题面指定的 commit \code{1ce47ea}。题面给出的 CUTLASS
pin 是短 hash \code{5c149f5}；构建时解析并 checkout 的完整 hash 为：

\begin{lstlisting}
5c149f52a436782210263fb2f19b354443a61c6a
\end{lstlisting}

\subsection{先阅读、后构建：关键源码事实}

读取顺序是 \code{TASK.md}、\code{BENCHMARK\_GB200.md}、
\code{docs/20260420-flashkda-v1-deep-dive.md}，再进入
\code{csrc/smxx/}。实际定位命令：

\begin{lstlisting}
cd assignment02/team/c1_flashkda
rg -n "CHUNK|kK2Threads|SM80_16x8x16|StateFP32" \
  FlashKDA/csrc/smxx
sed -n '1,220p' FlashKDA/docs/20260420-flashkda-v1-deep-dive.md
\end{lstlisting}

读到的事实如下。

\begin{enumerate}
  \item \code{fwd\_launch.cu} 中是 \code{constexpr int CHUNK = 16}；
        K1 grid 是 \code{(total\_tiles,H)}，K2 grid 是 \code{(N,H)}。
  \item \code{utils.cuh} 使用
        \code{SM80\_16x8x16\_F32BF16BF16F32\_TN} 和
        \code{SM80\_16x8x16\_F16F16F16F16\_TN}。
  \item K1 是 token-parallel：gate、L2 normalize、decay、\(L/M_{qk}\)
        与 \(16\times16\) Neumann inverse；K2 是 head-parallel：
        chunk recurrence、输出投影和最终 state。
  \item K2 有 192 threads：四个 MMA warp、一个 load warp、一个 store warp。
        每个 MMA warp 处理两个 \(16\times16\) column block。
  \item \code{SharedStorageK2::state\_acc} 是 bf16。所谓 fp32 state
        分支将 global fp32 转到 bf16 shared buffer，递推后再转回 fp32。
\end{enumerate}

\begin{decision}{源码层面的第一轮判断}
当前实现是围绕 \(16\times16\times16\) SM80 atom、bf16 shared state、
K1/K2 分离和 chunk 状态依赖共同设计的整体。后续替换 tcgen05 不是改一个
类型别名，而是可能触及 tile、shared layout、accumulator 生命周期和 CTA
划分的系统改动。
\end{decision}

\section{集群接入与环境确认}

\subsection{SSH 和 Slurm 使用约定}

实验使用两个 SSH alias：

\begin{lstlisting}
ssh infra5090
ssh infrab300
\end{lstlisting}

二者均是 Slurm 登录节点。登录节点只做 source、日志和调度查询；所有 CUDA
探测、编译、NCU、PyTorch CUDA 和 benchmark 都进入单卡 allocation。首轮的
资源探测形式如下：

\begin{lstlisting}
ssh infrab300 'hostname; sinfo -o "%P %D %G %N"'
ssh infrab300 '
  srun --partition=gpu --gres=gpu:1 --time=00:10:00 \
       bash -lc "hostname; nvidia-smi; nvcc --version"
'

ssh infra5090 'hostname; sinfo -o "%P %D %G %N"'
ssh infra5090 '
  srun --partition=lcpu-infra --gres=gpu:1 --time=00:10:00 \
       bash -lc "hostname; nvidia-smi; nvcc --version"
'
\end{lstlisting}

\subsection{确认到的硬件和工具链}

\begin{longtable}{@{}L{0.27\linewidth}L{0.31\linewidth}L{0.31\linewidth}@{}}
\toprule
项目 & \code{infrab300} & \code{infra5090} \\
\midrule
\endhead
\bottomrule
\endlastfoot
登录节点 & \code{b300-login} & \code{slurm-login} \\
partition & \code{gpu} & \code{lcpu-infra} \\
实际 GPU 节点 & \code{dev-slurm} & \code{gj-5090-1} \\
GPU & NVIDIA B300 SXM6 AC & NVIDIA GeForce RTX 5090 \\
Compute capability & 10.3 & 12.0 \\
Driver & 580.126.09 & 580.142 \\
CUDA toolkit & 13.0 / 13.0.88 & 13.0 \\
FlashKDA arch & \code{103a} & \code{120a} \\
\end{longtable}

\begin{record}{环境中的两个实际坑}
第一，B300 登录节点的 \code{nvidia-smi} 不在 PATH；系统 Python 也没有
\code{torch}。进入 \code{srun} 并使用 assignment venv 后，
\code{torch.cuda.is\_available()} 才为真。第二，不能在 5090 登录节点上
直接承担长时间 clone/compile；资源重操作全部改为调度节点或复用已验证的
B300 文件树。
\end{record}

\section{Python 依赖、CUTLASS 和 extension 构建}

\subsection{统一的远端目录变量}

\begin{lstlisting}
export C1_REPO=/home/lcpu/<user>/<repo>
export C1_A2=$C1_REPO/assignment02
export C1_DIR=$C1_A2/team/c1_flashkda
export C1_FLASH=$C1_DIR/FlashKDA
export C1_PY=$C1_A2/.venv/bin/python
\end{lstlisting}

B300 venv 为 Python 3.12，5090 为 Python 3.10.12；两边均使用
\code{torch 2.13.0+cu130} 和 \code{triton 3.7.1}。系统 Python 不能替换
这个 venv。

\subsection{依赖安装及其触发原因}

先安装 FLA reference runtime；第一次运行官方测试后发现缺 Ninja，继续运行
又发现测试要生成图像而缺 Matplotlib。因此两台机器在同一 venv 补齐：

\begin{lstlisting}
cd "$C1_A2"
uv pip install --python "$C1_PY" "flash-linear-attention>=0.5.0"
uv pip install --python "$C1_PY" ninja==1.13.2 matplotlib==3.11.1
\end{lstlisting}

得到 \code{flash-linear-attention 0.5.2}。Ninja 服务于 FLA/Triton 的
JIT helper；Matplotlib 是 \code{tests/test\_fwd.py} 输出
\code{plot.png}/\code{plot\_varlen.png} 的真实依赖，不是为了美化报告而装。

\subsection{B300：CUTLASS checkout}

\FlashKDA{} 快照没有题目指定的 CUTLASS submodule。B300 的构建流程：

\begin{lstlisting}
cd "$C1_FLASH"
git clone --filter=blob:none --no-tags --no-checkout \
  https://github.com/NVIDIA/cutlass.git cutlass
cd cutlass
git checkout --detach 5c149f52a436782210263fb2f19b354443a61c6a
\end{lstlisting}

直接以短 hash fetch 时 remote 未提供对应 ref；换成完整 commit 后成功。
最终 B300 CUTLASS 工作树约 191\,MiB，包含
\code{include}、\code{examples/common} 和 \code{tools/util/include}。

\subsection{5090：clone 无进展后的源码传输}

5090 登录节点的直接 CUTLASS clone 长时间无进展。没有删除不完整目录，也
没有冒充它可用；改用已经验证的 B300 checkout，通过 SSH 流式 tar 传输，显式
排除 Git metadata：

\begin{lstlisting}
ssh infrab300 "
  tar --exclude=cutlass/.git -C \"$C1_FLASH\" -czf - cutlass
" | ssh infra5090 "
  tar -xzf - -C \"$C1_FLASH\"
"
\end{lstlisting}

传输后检查 \code{include/cutlass/cutlass.h}、\code{examples/common}、
\code{tools/util/include}。目标树约 155\,MiB；无 \code{.git} 不影响
\code{setup.py} 所需 include 路径。

\subsection{两种 arch 的 editable build}

\begin{lstlisting}
# B300
srun --partition=gpu --gres=gpu:1 --time=00:30:00 \
  --chdir="$C1_FLASH" bash -lc '
    export FLASH_KDA_CUDA_ARCHS=103a
    export MAX_JOBS=8
    uv pip install --python "$C1_PY" --no-build-isolation --editable .
  '

# RTX 5090
srun --partition=lcpu-infra --gres=gpu:1 --time=00:30:00 \
  --chdir="$C1_FLASH" bash -lc '
    export FLASH_KDA_CUDA_ARCHS=120a
    export MAX_JOBS=8
    uv pip install --python "$C1_PY" --no-build-isolation --editable .
  '
\end{lstlisting}

两次 build 均成功，Python 包显示 \code{0.0.1+86106d7}。生成的动态库分别
是 \code{flash\_kda\_C.cpython-312-...so}（B300）和
\code{flash\_kda\_C.cpython-310-...so}（5090）。

\section{官方正确性复现}

\subsection{测试覆盖和命令}

\code{tests/test\_fwd.py} 覆盖 fixed
\(B=1,T=8192,H=96,D=128\) 的 PyTorch exact equality、varlen
\([1300,547,2048,963,271,3063]\) 的 exact equality，以及多种 gate/bias
分布下与 FLA/Triton 的容差对拍。

\begin{lstlisting}
cd "$C1_FLASH"
"$C1_PY" -u tests/test_fwd.py \
  | tee "$C1_DIR/results/test_fwd_<machine>.log"
\end{lstlisting}

实际执行时该命令位于 \code{srun --gres=gpu:1} allocation 内；这里把
Slurm 外壳省略，仅显示可复跑的程序主体。

\subsection{B300：通过}

\begin{longtable}{@{}L{0.34\linewidth}L{0.54\linewidth}@{}}
\toprule
检查项 & 记录结果 \\
\midrule
\endhead
\bottomrule
\endlastfoot
fixed vs PyTorch & avg/max rtol = 0；avg/max atol = 0；exact match \\
varlen vs PyTorch & avg/max rtol = 0；avg/max atol = 0；exact match \\
FLA fixed 10 个 case & FlashKDA output ratio 约 \(2.92\times10^{-3}\) 到 \(4.32\times10^{-3}\)，state ratio 约 \(3.43\times10^{-3}\) 到 \(9.63\times10^{-3}\) \\
FLA varlen 10 个 case & FlashKDA output ratio 约 \(3.03\times10^{-3}\) 到 \(5.15\times10^{-3}\)，state ratio 约 \(3.35\times10^{-3}\) 到 \(8.93\times10^{-3}\) \\
脚本终态 & \code{Assert results: Success} \\
\end{longtable}

应严格区分：PyTorch reference 是逐元素相等；FLA/Triton 是不同数值路径，
测试接受有限 ratio 并可能给 warning，最终仍通过。

\subsection{5090：严格 equality 未通过}

5090 第一个 fixed exact assertion 即停止：

\begin{lstlisting}
output | avg_rtol: 4.340156678495077e-09
output | max_rtol: 1.3563368383984198e-06
output | avg_atol: 7.634318066607193e-09
output | max_atol: 0.0625
AssertionError: output mismatch between kernel and torch ref
\end{lstlisting}

本轮只记录事实：5090 已构建并能 benchmark，但未通过
\code{torch.equal}。上述数值很小，可能涉及架构相关舍入/FMA 顺序；然而
allclose 分布、final state 和重复运行还未做，所以不能把原因写成定论。

\section{性能基线：方法、日志和数字}

\subsection{计时协议}

\code{benchmarks/bench\_fwd.py} 为每个调用创建 CUDA event pair，在 warmup
后收集 mean/min/max。比较五种路径：FlashKDA bf16 state、无 state、fp32
global state interface，以及 FLA \code{chunk\_kda} 和
\code{chunk\_gated\_delta\_rule}。

官方 H=96 形状使用题面配置：

\begin{lstlisting}
cd "$C1_FLASH"
"$C1_PY" -u benchmarks/bench_fwd.py \
  --mode all --H 96 --D 128 \
  --warmup 30 --iters 200 --repeats 5 \
  | tee "$C1_DIR/results/bench_official_h96_<machine>.log"
\end{lstlisting}

H=64 是第一轮早期的附加实验，参数为 \code{10/50/3}；其功能是暴露 shape
敏感性，不作为与 H=96 官方组等权的统计比较。

\subsection{H=96 官方参数结果}

单位是 ms，FlashKDA 使用 bf16 state。加速比定义为
\(\text{reference mean}/\text{FlashKDA mean}\)。

\begin{longtable}{@{}L{0.18\linewidth}rrrrr@{}}
\toprule
机器 / case & FlashKDA & chunk\_kda & 对 chunk & gdn & 对 gdn \\
\midrule
\endhead
\bottomrule
\endlastfoot
B300 fixed & 1.0324 & 2.3723 & 2.298x & 1.3056 & 1.265x \\
B300 varlen mixed & 0.8591 & 2.3905 & 2.783x & 1.3120 & 1.527x \\
B300 varlen 1024x8 & 0.7003 & 2.3482 & 3.353x & 1.2618 & 1.802x \\
5090 fixed & 2.6148 & 5.4165 & 2.071x & 3.0850 & 1.180x \\
5090 varlen mixed & 2.3280 & 5.4629 & 2.347x & 3.1463 & 1.352x \\
5090 varlen 1024x8 & 2.0546 & 5.4488 & 2.652x & 3.1189 & 1.518x \\
\end{longtable}

\subsection{H=64 附加形状}

\begin{longtable}{@{}L{0.18\linewidth}rrrrr@{}}
\toprule
机器 / case & FlashKDA & chunk\_kda & 对 chunk & gdn & 对 gdn \\
\midrule
\endhead
\bottomrule
\endlastfoot
B300 fixed & 1.6077 & 2.4776 & 1.541x & 1.3337 & 0.830x \\
B300 varlen mixed & 1.1159 & 2.5820 & 2.314x & 1.4329 & 1.284x \\
B300 varlen 1024x8 & 0.8034 & 2.3856 & 2.969x & 1.2686 & 1.579x \\
5090 fixed & 2.2103 & 3.5283 & 1.596x & 2.0090 & 0.909x \\
5090 varlen mixed & 1.6745 & 3.5726 & 2.134x & 2.0698 & 1.236x \\
5090 varlen 1024x8 & 1.4718 & 3.6000 & 2.446x & 2.0631 & 1.402x \\
\end{longtable}

\begin{decision}{性能的第一轮结论}
FlashKDA 在 H=96 的三类官方形状上对 FLA chunk\_kda 稳定领先，规则
1024x8 varlen 最强。它并不在所有 shape 上胜过 gdn：B300 H=64 fixed 的
gdn 为 1.3337\,ms，快于 FlashKDA 的 1.6077\,ms。报告必须保留这种反例。
\end{decision}

\section{NCU 和 SASS：指令及资源证据}

\subsection{NCU 运行、脚本问题和修复}

官方 \code{benchmarks/ncu.sh} 是有效模板，但使用 bare \code{pip}；
uv venv 中没有该命令。保留其 kernel filter 与 profile 意图，直接调用 venv
Python：

\begin{lstlisting}
cd "$C1_FLASH"
ncu --set full --kernel-name-base function \
  -k "regex:_flash_kda_fwd_(prepare|recurrence)" \
  --clock-control none --import-source yes --source-folders . \
  --export "$C1_DIR/results/ncu_fixed_h96_b300.ncu-rep" \
  "$C1_PY" -u benchmarks/bench_fwd.py \
    --mode fixed --H 96 --D 128 --warmup 0 --iters 5 --repeats 1
\end{lstlisting}

NCU 警告 6 个 CTC cross-chip metrics 无法访问；其余数据有效。profile 的
时钟和 instrumentation 会污染时间，所以不拿其 duration 代替 event benchmark。

\subsection{K1/K2 代表性指标}

\begin{longtable}{@{}L{0.27\linewidth}L{0.29\linewidth}L{0.30\linewidth}@{}}
\toprule
指标 & K1 prepare & K2 recurrence \\
\midrule
\endhead
\bottomrule
\endlastfoot
Block / grid & 256 / 49152 & 192 / 96 \\
Dynamic smem & 21.25 KiB & 98.432 KiB；allocated 约 99.456 KiB \\
Registers/thread & --- & 73--74 \\
Theoretical occupancy & --- & 18.75\% \\
Achieved occupancy & 97.07\% & 约 9.37--9.38\% \\
Compute throughput & 约 73.56\% & 约 20.8--21.0\% \\
Memory throughput & 约 71.32\% & 约 38\% \\
DRAM throughput & 约 34.95\% & 约 10.7--11.2\% \\
Tensor pipe active & --- & 约 19.5--20.25\% \\
\end{longtable}

K1 有大量 tile 并行，资源利用高。K2 fixed H=96 的 grid 只有 96 CTA，
且每 CTA 接近 100\,KiB shared state/pipeline；它呈现低 tensor 活跃度、
低 DRAM 吞吐和状态依赖，不宜在第一轮草率贴上唯一的 memory-bound 或
compute-bound 标签。

\subsection{SASS}

\begin{lstlisting}
cuobjdump --dump-sass flash_kda_C.cpython-312-x86_64-linux-gnu.so \
  > "$C1_DIR/results/sass/flash_kda_C.sm103.sass"
cuobjdump --dump-sass flash_kda_C.cpython-310-x86_64-linux-gnu.so \
  > "$C1_DIR/results/sass/flash_kda_C.sm120.sass"

grep -E -c 'HMMA\.16816\.F16' results/sass/*.sass
grep -E -c 'HMMA\.16816\.F32\.BF16' results/sass/*.sass
grep -E -c 'WGMMA|TCGEN05' results/sass/*.sass
\end{lstlisting}

\begin{center}\small
\begin{tabular}{@{}lrr@{}}
\toprule
pattern & B300 SM103 & 5090 SM120 \\
\midrule
\code{HMMA.16816.F16} & 24 & 24 \\
\code{HMMA.16816.F32.BF16} & 1520 & 1520 \\
\code{WGMMA} & 0 & 0 \\
\code{TCGEN05} & 0 & 0 \\
\bottomrule
\end{tabular}
\end{center}

所以「在 Blackwell 上运行」不等于「用 Blackwell 新 MMA」。C1 题面中
SM80 路径仍在 Blackwell 上执行的前提，被源码和两份 SASS 同时支持。

\section{SM100 challenge：tcgen05 的正例、负例和边界}

\subsection{纸面约束}

CUTLASS \code{cute/arch/mma\_sm100\_umma.hpp} 中
\code{SM100\_MMA\_F16BF16\_SS} 有如下编译期限制：

\begin{lstlisting}
static_assert(M == 64 || M == 128,
  "SM100_MMA_F16BF16 M-mode size should be 64 or 128 ...");
static_assert((N % 8 == 0) && (8 <= N) && (N <= 256), ...);
\end{lstlisting}

生产 K2 的局部块却是 \(16\times16\times16\)，并采用 warp-local register
accumulator。tcgen05 从 shared memory（或 TMEM）读取操作数、把 accumulator
写入 TMEM，额外需要 TMEM allocator、completion barrier 和 TMEM-to-register
epilogue。由此可知，直接替换 atom 名称无法保留现有数据路径。

\subsection{可运行的 B300/SM103a 正例}

新增 \code{challenge/tcgen05\_sm103\_gemm.cu}，来源是 pin 后 CUTLASS 的
\code{examples/cute/tutorial/blackwell/01\_mma\_sm100.cu}，保留 SPDX 头。
原 tutorial host check 只允许 10.0/10.1；B300 为 10.3，故只放宽 host check，
仍以 \code{sm\_103a} 编译。

\begin{lstlisting}
cd "$C1_DIR/challenge"
export CUTLASS_ROOT="$C1_FLASH/cutlass"
./build_sm103.sh 128 256 64

nvcc -std=c++17 -O3 -arch=sm_103a --expt-relaxed-constexpr \
  -I"$CUTLASS_ROOT/include" \
  -I"$CUTLASS_ROOT/tools/util/include" \
  -I"$CUTLASS_ROOT/examples/common" \
  -I. -o tcgen05_sm103 tcgen05_sm103_gemm.cu
\end{lstlisting}

该 GEMM 使用 F16xF16$\rightarrow$F32 的
\code{SM100\_MMA\_F16BF16\_SS<...,128,256>}；四个 K=16 指令组成
\(128\times256\times64\) 工作例。第一次误把运行时 K 设为 16，tutorial
提示 problem shape 不被 K=64 tiler 整除；改为 \code{128 256 64} 后，CPU
reference infinity-norm 差为 0，relative error 为 0，并打印
\code{Execution is successful.}。

\subsection{故意构造的 16x16 负例}

在临时副本中把 atom 设为 \(M=N=16\)，同一 \code{nvcc -arch=sm\_103a}
命令立即失败：

\begin{lstlisting}
SM100_MMA_F16BF16 M-mode size should be 64 or 128
UMMA_1SM M-mode size should be 64 or 128
\end{lstlisting}

完整输出见 \code{results/tcgen05\_m16\_compile.log}。这是 trait 在实例化时
拒绝 M=16，而非运行慢、参数差或运行时不稳定。

\subsection{轻量 tcgen05 NCU}

\begin{lstlisting}
ncu --target-processes all --kernel-name-base function \
  -k "regex:gemm_device" \
  --metrics gpu__time_duration.avg,\
sm__pipe_tensor_cycles_active.avg.pct_of_peak_sustained_elapsed,\
sm__warps_active.avg.pct_of_peak_sustained_elapsed,\
dram__bytes_read.sum,dram__bytes_write.sum \
  --csv ./tcgen05_sm103 128 256 64
\end{lstlisting}

CSV 显示 block=128、grid=\((1,1,1)\)，kernel 名含
\code{SM100\_MMA\_F16BF16\_SS<...,128,256,...>}。profile duration 约
25.1\,ms，旧通用 tensor metric 为 0；二者都不用于同 FlashKDA 做性能结论，
因为这是单 CTA 小 tutorial，NCU instrumentation 又会显著放大时间。该实验
的职责是证明 tcgen05 能在 SM103a 正确执行，并证明 16x16 不是可接受 tile。

\begin{decision}{只换指令路线的首轮结论}
局部 patch 不可行。SM100 v2 至少需要重做 K2 M/N tile、head/CTA 划分、
shared swizzle/descriptor、TMEM lifetime、barrier 和 epilogue。负结果支持
v1 保持 SM80 通用路径；但没有证明完整 SM100 v2 必然不值得做。
\end{decision}

\section{CHUNK 数值范围和 state 精度 probe}

\subsection{CHUNK range probe}

K1 采用 base-2 gate：
\[
g_{\log_2}=\texttt{lower\_bound}\log_2(e)\,
\sigma(\exp(A_{\log})(g+\texttt{dt\_bias})),\qquad
\texttt{lower\_bound}=-5.
\]
对一个 chunk 做 \(\exp_2(\mathrm{cumsum}(g_{\log_2}))\) 后写回 bf16。
bf16 最小 normal/subnormal 约为 \(2^{-126}\)/\(2^{-133}\)，CHUNK 变大时
累计负指数更容易下溢。

新增 \code{challenge/chunk\_range\_probe.py}，用 1,048,576 个独立 scalar，
seed 0、B300 CUDA、benchmark 同型分布（bf16 normal g，float32 uniform
dt\_bias/A\_log）执行：

\begin{lstlisting}
"$C1_PY" challenge/chunk_range_probe.py \
  --samples $((1<<20)) --seed 0 \
  | tee "$C1_DIR/results/chunk_range_b300.log"
\end{lstlisting}

\begin{center}\small
\begin{tabular}{@{}rrrr@{}}
\toprule
CHUNK & log2 最小值 & exp2 最小值 & bf16 zero fraction \\
\midrule
16 & -103.066 & \(9.419\times10^{-32}\) & 0.000000 \\
32 & -189.270 & 0 & 0.098475 \\
64 & -347.805 & 0 & 0.547464 \\
\bottomrule
\end{tabular}
\end{center}

这是标量范围 probe，不是完整 kernel 正确性。它已经说明在当前 lower bound
下，CHUNK=32 开始存在明显 bf16 zero，CHUNK=64 更严重。大 CHUNK 必须先设计
rescale 或更高精度存储；不能只改 \code{constexpr CHUNK=16}。同时 inverse
矩阵元素数从 \(16^2=256\) 增至 \(32^2=1024\)、\(64^2=4096\)，现有
Neumann helper、workspace、layout 和 SM80 分块都会同步重写。

\subsection{state precision probe}

源码表明 fp32 global state 会转为 bf16 \code{state\_acc}，结束时再转回
fp32。因此新增 \code{challenge/state\_precision\_probe.py}，固定
\(B=1,T=1024,H=4,D=128\)、seed 0，比较：

\begin{enumerate}
  \item FlashKDA bf16 global state；
  \item FlashKDA fp32 global state interface；
  \item FLA chunk\_kda fp32 reference。
\end{enumerate}

\begin{lstlisting}
"$C1_PY" challenge/state_precision_probe.py \
  | tee "$C1_DIR/results/state_precision_b300.log"
\end{lstlisting}

\begin{center}\small
\begin{tabular}{@{}lrr@{}}
\toprule
比较 & max abs & mean abs \\
\midrule
bf16-state vs fp32-state output & 0 & 0 \\
bf16-state vs fp32-state final & 0 & 0 \\
FlashKDA output vs FLA & \(4.882812\times10^{-4}\) & \(1.289556\times10^{-5}\) \\
FlashKDA final state vs FLA & \(3.915220\times10^{-3}\) & \(9.135994\times10^{-5}\) \\
\bottomrule
\end{tabular}
\end{center}

near-zero 元素使 max relative error 很大（约 \(7.07\times10^1\) 和
\(9.99\times10^1\)），故本文采用 absolute max/mean 描述整体差异。

\begin{decision}{state 的首轮结论}
bf16/fp32 interface 的逐元素相同并不表示 fp32 internal state 没价值；
它只表示当前两条接口在 bf16 \code{state\_acc} 汇合。若第二轮要研究
内部 fp32 state，必须实现真正的 fp32 shared recurrence 或独立模拟它。
\end{decision}

\section{失败尝试和修复记录}

\begin{longtable}{@{}L{0.23\linewidth}L{0.31\linewidth}L{0.35\linewidth}@{}}
\toprule
尝试/现象 & 观察 & 修复或保留结论 \\
\midrule
\endhead
\bottomrule
\endlastfoot
登录节点探 GPU/import torch & B300 login 无 nvidia-smi PATH，系统 Python 无 torch & GPU/Python CUDA 均进入 srun，显式使用 assignment venv \\
首次官方测试 & 缺 Ninja，FLA helper/JIT 无法开始 & 安装 \code{ninja==1.13.2} \\
测试继续执行 & 缺 Matplotlib，绘图步骤失败 & 安装 \code{matplotlib==3.11.1} \\
B300 CUTLASS 短 hash & remote 不接受短 hash ref & checkout 完整 hash \\
5090 直接 clone & 登录节点 clone 无进展 & 保留不完整目录；从 B300 传输已验证源码 \\
官方 ncu.sh & 假定 bare pip，uv 环境没有 & 直接以 venv Python 执行 NCU \\
tcgen tutorial K=16 & K=64 tiler 不整除，拒绝 OOB & 改为 128,256,64；非 tcgen 指令失败 \\
tcgen M=16 & 编译期 static assertion & 用作 tile 不可直接替换的证据 \\
5090 torch.equal & 微小非零误差导致 strict fail & 不伪称通过；第二轮做 allclose/final-state \\
\end{longtable}

\section{版本控制、产物和最短复现}

第一轮已推送的提交是：

\begin{lstlisting}
3e0e391 assignment02: add C1 baseline and SM100 challenge evidence
\end{lstlisting}

提交含报告草稿、challenge 源码、构建脚本和轻量日志；未包含
\code{ncu\_fixed\_h96\_b300.ncu-rep}（约 90\,MiB）、完整 SASS dump 以及
任何与 C1 无关文件。

在 B300 上的最短复现入口：

\begin{lstlisting}
ssh infrab300
export C1_REPO=/home/lcpu/<user>/<repo>
export C1_A2=$C1_REPO/assignment02
export C1_DIR=$C1_A2/team/c1_flashkda
export C1_FLASH=$C1_DIR/FlashKDA
export C1_PY=$C1_A2/.venv/bin/python

srun --partition=gpu --gres=gpu:1 --time=00:30:00 \
  --chdir="$C1_FLASH" "$C1_PY" -u tests/test_fwd.py
srun --partition=gpu --gres=gpu:1 --time=00:30:00 \
  --chdir="$C1_FLASH" "$C1_PY" -u benchmarks/bench_fwd.py \
  --mode all --H 96 --D 128 --warmup 30 --iters 200 --repeats 5
srun --partition=gpu --gres=gpu:1 --time=00:20:00 \
  --chdir="$C1_DIR/challenge" bash -lc '
    export CUTLASS_ROOT="$C1_FLASH/cutlass"
    ./build_sm103.sh 128 256 64
    "$C1_PY" chunk_range_probe.py --samples $((1<<20)) --seed 0
    "$C1_PY" state_precision_probe.py
  '
\end{lstlisting}

\section{第一轮里程碑：结论强度与第二轮入口}

\subsection{已有充分证据的结论}

\begin{enumerate}
  \item B300 可构建并通过官方 PyTorch fixed/varlen exact 对拍。
  \item B300/5090 生产 extension 的实际 MMA 仍是 HMMA/SM80，而非 WGMMA/TCGEN05。
  \item H=96 官方形状对 FLA chunk\_kda 有 2.07x--3.35x 加速，但优势依赖 shape。
  \item tcgen05 的最小 M tile 与 TMEM accumulator 模型和 K2 的 16x16 SM80
        设计直接冲突。
  \item 当前 lower bound 下 CHUNK=32 已存在 bf16 exponent 下溢，64 更严重。
  \item 当前 fp32 state 只改变 global I/O，不改变 bf16 internal recurrence。
\end{enumerate}

\subsection{还不能下结论的事项}

\begin{enumerate}
  \item tutorial microbench 不能证明完整 SM100 v2 必然更慢；
  \item 5090 strict mismatch 尚不能确定为舍入顺序；
  \item 标量 range probe 不能替代完整 CHUNK=32/64 kernel 的精度和性能测量；
  \item K2 低 occupancy 不能被简化成唯一的 bandwidth-bound 或 compute-bound；
  \item 接口 fp32 的无差异不能外推为 internal fp32 没有价值。
\end{enumerate}

\subsection{第二轮建议顺序}

\begin{longtable}{@{}L{0.12\linewidth}L{0.34\linewidth}L{0.39\linewidth}@{}}
\toprule
优先级 & 工作 & 可验收交付 \\
\midrule
\endhead
\bottomrule
\endlastfoot
1 & 5090 allclose/final-state 诊断 & 固定 seed、重复运行、输出和 state 误差分布 \\
2 & K2 并行度纸面上界 & 多 head/CTA、persistent、2-CTA 的 smem/reg/grid/依赖反例 \\
3 & 选择一个并行度 prototype & reference 对拍、occupancy、端到端时间、负结果机制 \\
4 & 大 CHUNK reference/rescale & inverse error、下溢、rescale、资源和速度逐项记录 \\
5 & v2 SM100 专版决策 & 将可移植性、工程成本、峰值潜力组织成最终结论 \\
\end{longtable}

\begin{record}{第二轮实验纪律}
每一个新数字前都记录 GPU、commit、输入 shape、seed、state dtype、
warmup/iters/repeats 和是否处于 profiler。C1 的价值不在于写出一条
「用了 tcgen05」的结论，而在于把 tile、数值范围、状态依赖和实测指标连接为
可复查、可反驳的因果链。
\end{record}

\clearpage
\section{第二轮：精确性诊断与 K2 并行度方案}

本章从第二轮开始，按时间追加；在本章开头先写下实验意图，避免见到结果后再
倒推问题。第二轮的第一件事不是试图优化 kernel，而是界定 5090 上的
\code{torch.equal} 失败究竟是非确定性、跨 chunk 才出现的舍入差异、state
接口差异，还是更实质的错误。若不先完成这个区分，任何 K2 并行改写都没有
可信的正确性基线。

\subsection{R2.1：5090 精确性诊断---预注册设计}

第一轮在 5090 上的官方 fixed test 报告极小但非零差异；B300 则 exact match。
这只说明 ``bitwise identical'' 不成立，尚不足以说明算法错。新增独立脚本
\code{challenge/diagnose\_5090\_exactness.py}，不改动 \FlashKDA{} 源码，
直接复用官方 \code{tests/torch\_ref.py} 与官方 fixed 输入构造。它对每个 case
记录下列项目：

\begin{enumerate}
  \item 相同 reference 连续运行两次是否 exact，排除 reference 自己的非确定性；
  \item 相同 kernel 连续运行至少三次是否 exact，排除 kernel 执行非确定性；
  \item output 和 final state 分别相对 reference 的 mismatch 个数、比例、
        mean/max absolute error；
  \item bf16 输出的 mismatch 是否超过 reference 相邻 bf16 表示值的一步；
  \item 至多四个 mismatch 的 index、两端值和绝对误差，提供可人工复核的样本；
  \item device name/capability、Torch/CUDA、extension 路径和 TF32 设置。
\end{enumerate}

\paragraph{边界矩阵。} \code{boundary} preset 覆盖 \(T=16,17,64,1024\)、
\((T,H)=(8192,1),(8192,96)\) 及 \(T=8192,H=96\) 的 bf16/fp32 state 接口。
\(T=16\) 恰好一个 chunk，\(T=17\) 第一次触发跨 chunk recurrence；H=1 与
H=96 区分单个 K2 CTA 的算术问题和多 CTA 数量/调度因素。fp32 case 只检验
公开 state 接口，不误称其已经检验了内部 fp32 recurrence。

\paragraph{实际执行命令、脚本身份和作业。} 两台机器上运行的是同一份脚本；
复制后分别以 \code{sha256sum} 核验，SHA-256 均为
\code{ba6ba8b977c21336697f1b29bf75daf662089886ada9e73f9f9e567361b65566}。
诊断不是 benchmark，wall clock 不用于任何性能判断。为遵守登录节点约定，
Python、JIT 与 CUDA 均在 \code{srun} 分配内发生。5090 远端主仓库停在旧
commit 且 \code{assignment02/} 是未跟踪目录；因此没有 \code{git pull}、reset
或清理远端工作树，只复制这个已经提交的独立脚本到 C1 的 \code{challenge/}。

\begin{lstlisting}
# 5090：job 9875，partition=lcpu-infra，node=gj-5090-1，最终 COMPLETED (3m11s)
ssh infra5090
export C1_REPO=/home/lcpu/<user>/<repo>
export C1_A2=$C1_REPO/assignment02
export C1_DIR=$C1_A2/team/c1_flashkda
export C1_FLASH=$C1_DIR/FlashKDA
export C1_PY=$C1_A2/.venv/bin/python

cd "$C1_FLASH"
srun --partition=lcpu-infra --gres=gpu:1 --time=00:30:00 --chdir="$PWD" \
  "$C1_PY" ../challenge/diagnose_5090_exactness.py \
    --preset boundary --seed 0 --repeats 3 \
    --json ../results/diagnose_5090_boundary.json \
    2>&1 | tee ../results/diagnose_5090_boundary.log

# B300：第一次 job 23441 失败；第二次 job 23442 成功。
# 失败/修复原因见本节的环境记录；成功运行显式把 venv/bin 加入 PATH。
ssh infrab300
export C1_REPO=/home/lcpu/<user>/<repo>
export C1_A2=$C1_REPO/assignment02
export C1_DIR=$C1_A2/team/c1_flashkda
export C1_FLASH=$C1_DIR/FlashKDA
export C1_PY=$C1_A2/.venv/bin/python
cd "$C1_FLASH"
srun --partition=gpu --gres=gpu:1 --time=00:30:00 --chdir="$PWD" \
  env PATH="$C1_A2/.venv/bin:$PATH" "$C1_PY" -u \
    ../challenge/diagnose_5090_exactness.py \
    --preset boundary --seed 0 --repeats 3 \
    --json ../results/diagnose_b300_boundary_path.json \
    2>&1 | tee ../results/diagnose_b300_boundary_path.log
\end{lstlisting}

\subsection{R2.1：执行记录与结果}

\subsubsection{环境与一次可恢复的启动失败}

5090 job 9875 成功结束。JSON 中记录的设备是 NVIDIA GeForce RTX 5090，
compute capability 12.0；Python 是 3.10.12，\Fla{} 相关环境为
Torch \code{2.13.0+cu130}/CUDA \code{13.0}，且
\code{torch.backends.cuda.matmul.allow\_tf32=False}。

B300 的首次 job 23441 在导入官方 \code{tests/torch\_ref.py} 时失败，
错误不是 FlashKDA kernel，而是 PyTorch \code{load\_inline} 提示缺 Ninja。
随后的只读检查证明 \code{\$C1\_A2/.venv/bin/ninja} 已存在，真正原因是以
绝对路径启动 Python 时没有 activate venv，故 venv 的 \code{bin/} 不在
\code{PATH}。该失败 stdout 被保留为
\code{results/diagnose\_b300\_boundary\_no\_venv\_path.log}；没有重装依赖。
第二次 job 23442 将 \code{\$C1\_A2/.venv/bin} 前置到 \code{PATH} 后成功。
B300 JSON 中设备为 NVIDIA B300 SXM6 AC、compute capability 10.3；Python
是 3.12.13，Torch/CUDA/TF32 设置与 5090 JSON 相同。

\subsubsection{位级稳定性和两卡对照}

所有 7 个 case 中，reference 连续执行两次均 exact；每个设备的 FlashKDA
连续执行三次也均输出/最终 state exact。这排除了本诊断条件下的随机调度或
未初始化输出。下表的 ``5090 out/state'' 写作
``mismatch 个数/总元素数（最大绝对误差）''；out 的 ``$>1$ ulp'' 是不止相邻
bf16 表示值一步的 mismatch 数。B300 每一格都是零，故压缩为 ``0（exact）''。

{\footnotesize
\setlength{\tabcolsep}{2pt}
\begin{longtable}{@{}L{0.17\linewidth}L{0.13\linewidth}L{0.23\linewidth}L{0.23\linewidth}L{0.10\linewidth}@{}}
\toprule
case & \((T,H,\mathrm{state})\) & 5090 output & 5090 final state & B300 \\
\midrule
\endhead
\bottomrule
\endlastfoot
T16-H96 & \((16,96,\mathrm{bf16})\) & 5656/196608 (0.03125; $>1$ ulp: 1957) & 1143/1572864 (0.001953125) & 0 (exact) \\
T17-H96 & \((17,96,\mathrm{bf16})\) & 5671/208896 (0.03125; $>1$ ulp: 1962) & 189/1572864 (0.00048828125) & 0 (exact) \\
T64-H96 & \((64,96,\mathrm{bf16})\) & 5701/786432 (0.03125; $>1$ ulp: 1967) & 6/1572864 (0.0000038147) & 0 (exact) \\
T1024-H96 & \((1024,96,\mathrm{bf16})\) & 5079/12582912 (0.0078125; $>1$ ulp: 1719) & 142/1572864 (0.0001220703) & 0 (exact) \\
T8192-H1 & \((8192,1,\mathrm{bf16})\) & 55/1048576 (0.0001220703; $>1$ ulp: 18) & 0/16384 (0) & 0 (exact) \\
T8192-H96 & \((8192,96,\mathrm{bf16})\) & 9110/100663296 (0.0625; $>1$ ulp: 3177) & 49/1572864 (0.0000076294) & 0 (exact) \\
T8192-H96 & \((8192,96,\mathrm{fp32})\) & 9110/100663296 (0.0625; $>1$ ulp: 3177) & 49/1572864 (0.0000076294) & 0 (exact) \\
\end{longtable}
}

对官方 \((T,H,D)=(8192,96,128)\) case，5090 output 的 mean absolute error
为 \(7.6343181\times10^{-9}\)，相对 reference mean-absolute-value 的比为
\(4.3401567\times10^{-9}\)。这说明总体均值很小，但不等于 ``只差一 ulp''：
该 case 有 3177 个 output mismatch 超过相邻 bf16 表示值的一步，最大值为
0.0625。因此报告同时保存计数、最大值和均值，不能只展示一个平均相对误差。

\begin{decision}{R2.1 的结论边界}
\textbf{证据支持：} 5090 的 strict mismatch 是可重复的；它在 \(T=16\) 的单
chunk case 已出现，所以不能归咎于跨 chunk recurrence；H=1 时仍有 55 个
output mismatch，因此也不能归咎于 ``多 head 才会出错''；bf16/fp32 的公开
state 接口使 H=96 的 output/state mismatch 计数完全相同，支持第一轮的源码
发现---它没有改变内部 bf16 recurrence。相同诊断在 B300 的所有 case 均 exact，
也复核了第一轮 B300 结论。

\textbf{证据尚不足以支持：} 这还不能证明 5090 kernel 算法错误，更不能从两张不同设备
推出唯一根因。两边 GPU capability 不同，Python 小版本也不同；尽管 Torch/CUDA
版本、脚本、seed、输入规则、TF32 设置和 extension API 一致，仍需下一步用
SASS/编译选项、reference 中 \code{load\_inline} 生成的 JIT kernel 以及更小
的单 MMA 对照去区分硬件/工具链/舍入路径。第二轮后续的 K2 prototype 正确性
采用 B300 exact 对拍为硬门槛；5090 只报告误差分布，不能使用 \code{torch.equal}
作为唯一通过标准。
\end{decision}

\begin{record}{可审计产物}
\code{results/diagnose\_5090\_boundary.log} 与 \code{.json} 是 job 9875 的
stdout/结构化摘要；\code{results/diagnose\_b300\_boundary\_path.log} 与
\code{.json} 是 job 23442 的成功记录；\code{diagnose\_b300\_boundary\_no\_venv\_path.log}
保留 job 23441 的 PATH 失败。这些均为小文件，已作为本轮证据进入 git。
\end{record}

\subsection{R2.2：K2 并行度的资源上界与候选方案筛选}

\subsubsection{问题重述：依赖在哪一维，空闲又在哪一维？}

K2 对一个 \((\text{sequence},\text{head})\) 的所有 chunk 按时间顺序递推：
第 \(t+1\) 个 chunk 必须读取第 \(t\) 个 chunk 更新后的 \(128\times128\)
state。因此不能把同一 head 的时间维 chunk 直接拆给多个 CTA 后独立运行。
但不同 head 互不读写同一 state，且 state 的 value-column（矩阵第二维）在
K2 的这条公式中也没有跨列归约：

\[
  S'_{:,j}=S_{:,j}\odot g + K_{\mathrm{restored}}^\mathsf{T}U_{:,j},
  \qquad
  O_{:,j}=Q_{\mathrm{decayed}}S_{:,j}+M_{qk}U_{:,j}.
\]

这里 \(j\) 是 value column。上述表达与源码的 Phase 1--6 一致：四个 MMA
warp 各处理两个连续 \(16\times16\) 列块；Phase 6 的
\code{for (bi=0; bi<2; ++bi)} 对各列块独立更新 \code{s\_acc}。末尾的
\code{compute\_barrier} 只是 CTA 内同步，不是跨列求和。这给 ``按 value 列
切成两个 CTA'' 提供了正确性的代数前提；它不等价于已经证明这种切分会更快。

\subsubsection{B300 资源测量与现有 K2 上界}

第一轮 NCU 对官方 \((T,H,D)=(8192,96,128)\) K2 记录的 launch 是
block=192 threads、grid=\((N,H)=(1,96)\)、dynamic shared memory=98,432 B，
registers/thread 为 73--74，theoretical occupancy=18.75\%，achieved
occupancy=9.37\%。为避免把 ``低 occupancy'' 误解为某个猜测，本轮又在 B300
job 23452 内读取设备属性；输出保存在
\code{results/k2\_resource\_limits\_b300.log}：148 SM、每 SM 2048 threads、
233,472 B shared memory、65,536 registers。第一次 job 23451 还尝试读取
PyTorch 不提供的 \code{regs\_per\_block} 属性，已打印前述有效字段后以
\code{AttributeError} 结束；该失败保留在实验记录中。

实际执行的成功命令等价于：

\begin{lstlisting}
ssh infrab300
export C1_A2=/home/lcpu/<user>/<repo>/assignment02
export C1_PY=$C1_A2/.venv/bin/python
srun --partition=gpu --gres=gpu:1 --time=00:05:00 \
  env PATH="$C1_A2/.venv/bin:$PATH" "$C1_PY" -c '
import torch
p = torch.cuda.get_device_properties(0)
for name in ("multi_processor_count", "max_threads_per_multi_processor",
             "shared_memory_per_multiprocessor", "regs_per_multiprocessor"):
    print(f"{name}={getattr(p, name, None)}")
'
\end{lstlisting}

由上述数值可得到 K2 在一个 SM 上的 resident-CTA 上界：
\[
\begin{aligned}
B_{\mathrm{smem}} &= \left\lfloor\frac{233472}{98432}\right\rfloor=2,\\
B_{\mathrm{thread}} &= \left\lfloor\frac{2048}{192}\right\rfloor=10,\\
B_{\mathrm{reg}} &\mathrel{\gtrsim}\left\lfloor\frac{65536}{192\times74}\right\rfloor=4,\\
B_{\mathrm{resident}} &\leq \min(2,10,4)=2.
\end{aligned}
\]

shared memory 是决定性项；两 CTA 需要 196,864 B，三 CTA 则 295,296 B，
超过 233,472 B。两 CTA 的 384 threads 对应
\(384/2048=18.75\%\)，与 NCU 的 theoretical occupancy 相符。相反，官方
unbatched grid 只有 96 CTA，小于 148 SM；调度器若优先摊到空闲 SM，就会让
active SM 上只有一个 192-thread CTA（\(192/2048=9.375\%\)），这与 NCU 的
9.37\% achieved 值一致。该对应关系是对调度现象的解释，不是由一条 metric
单独证明的因果结论。

\subsubsection{三个候选方案与反例}

\begin{longtable}{@{}L{0.19\linewidth}L{0.22\linewidth}L{0.25\linewidth}L{0.24\linewidth}@{}}
\toprule
方案 & 能改变的并行度 & 资源/正确性反例 & R2.2 判断 \\
\midrule
\endhead
\bottomrule
\endlastfoot
多 head/CTA & 把两个 head 顺序装进一个 CTA，grid 从 96 降到 48 & 每个 head 都有独立 state、TMA pipeline 和 512 个 chunk；顺序复用不能缩短临界路径，且更少 CTA 反而恶化 H=96 的 grid 不足 & 拒绝作为首个 prototype \\
persistent CTA & 用约 SM 数的 CTA 循环领取 head & H=96 小于 148 SM；148 个 persistent CTA 中至少 52 个一开始无 head 可领。H 很大时可能省 launch/调度开销，但不打破单 head 的时间依赖 & 留作 H\(\gg\)SM 的后续变体，不用于官方形状首测 \\
2 CTA/head，按 value column 切 64+64 & grid 从 96 提高到 192；每 CTA 负责半个 state、半个 v/out 列 & 列独立保证可表达；但两个 CTA 不能共享 smem，仍须各自 TMA 读取全宽的 \code{k\_decayed}、\code{q\_decayed}、\code{k\_restored}、\code{g\_total}、\code{INV} 和 \code{Mqk}，会重复全局流量和 TMA/同步开销 & 唯一值得实现的候选；预期不是必然加速 \\
\end{longtable}

\subsubsection{2-CTA 列切分的资源草算与 prototype 边界}

当前 K2 的可见 shared-memory 主体可按源码粗分为 state 32 KiB，三 stage input
约 \(3\times17{,}984\) B，两个 stage output 8 KiB，另有 pipeline/barrier
元数据；合计与 NCU 的 98,432 B 相符。若只把 value dimension 从 128 切为 64，
state 约降至 16 KiB，\code{v/out} 约各减半；但 q/k side 的三个
\(16\times128\) workspace tensor 不能减半。故每个列切 CTA 的乐观估计仍约
68--72 KiB，不是 49 KiB，更不是简单的 ``原来的一半''。两个 CTA/head 的
总 smem 也会高于现有一个 CTA/head；潜在收益来自增加 grid，而不是减少数据量。

实现该 prototype 必须同时改：\code{VOLayout} 与 output TMA 的列 offset、
\code{StateSmemLayout} 与 state load/store 的列 slice、四个 MMA warp 到两个
MMA warp 的映射、尾 tile 的手工 store、以及 \((N,H,2)\) 的 grid 映射。只在
launch 处把 grid.z 设为 2 会使两个 CTA 写同一 output/state，属于错误实验。
因此 R2.2 的结果是预先限定 R2.3 的改动边界和验收标准：B300 exact 对拍，
NCU 记录新 smem/regs/occupancy，event benchmark 与原 K2 的端到端时间比较，
并把重复 Q/K TMA 流量作为可能负结果的机制。

\begin{record}{R2.2 的结论强度}
已经量化了官方 H=96 K2 的 grid/SM 不匹配和 smem 的两-CTA 上界，也从源码验证
了列独立这一必要条件。尚未实现 2-CTA kernel，因而不能声称 occupancy 提升、
也不能声称会有端到端收益。R2.3 的代码改动必须是隔离的 compile-time prototype，
不能污染已验证的基线路径。
\end{record}

\clearpage
\section{第三轮：2-CTA/head 列切分 prototype}

\subsection{R3.1：reference-level 列切分 oracle}

第三轮的目标是把 R2.2 的列独立推导落实为可执行的正确性前置检查，随后才进入
CUDA implementation。新增 \code{challenge/vsplit\_reference\_probe.py}，它不
调用 FlashKDA extension，也不测量性能；它复用官方 \code{torch\_ref.py} 的
L2 normalization、gate、Neumann inverse、bf16 cast 及 fp32 FMA 数值顺序，
但将每个 head 的 value dimension 分为 \([0,64)\) 和 \([64,128)\) 两段。

每一段读取同一份全宽 \(Q/K\) side 中间量，但只读取和更新自己的
\code{state[V\_slice,K]}、\code{v[:,V\_slice]} 与 output columns。该 oracle
对未切分 \code{torch\_ref} 的 output/final state 做 exact 对拍。若 reference
层就无法 exact，则取消 CUDA 列切分；若通过，才说明把 state/output TMA layout
分列是一个数值上可继续实现的候选。

计划在 B300 的 GPU allocation 中运行下列矩阵：

\begin{lstlisting}
ssh infrab300
export C1_REPO=/home/lcpu/<user>/<repo>
export C1_A2=$C1_REPO/assignment02
export C1_DIR=$C1_A2/team/c1_flashkda
export C1_FLASH=$C1_DIR/FlashKDA
export C1_PY=$C1_A2/.venv/bin/python

srun --partition=gpu --gres=gpu:1 --time=00:20:00 --chdir="$C1_FLASH" \
  env PATH="$C1_A2/.venv/bin:$PATH" "$C1_PY" -u \
    ../challenge/vsplit_reference_probe.py --T 16 --split 64 --seed 0
srun --partition=gpu --gres=gpu:1 --time=00:20:00 --chdir="$C1_FLASH" \
  env PATH="$C1_A2/.venv/bin:$PATH" "$C1_PY" -u \
    ../challenge/vsplit_reference_probe.py --T 17 --split 64 --seed 0
srun --partition=gpu --gres=gpu:1 --time=00:20:00 --chdir="$C1_FLASH" \
  env PATH="$C1_A2/.venv/bin:$PATH" "$C1_PY" -u \
    ../challenge/vsplit_reference_probe.py --T 64 --split 64 --seed 0
\end{lstlisting}

\begin{record}{R3.1 验收条件}
三个长度必须同时满足 output 与 final state 的 \code{torch.equal}；\(T=16\)
覆盖单 chunk，\(T=17\) 覆盖首次跨 chunk state 传递，\(T=64\) 覆盖四次连续
state 更新。通过仅证明列切分的 reference 语义，不能代替 CUDA prototype 的
正确性、NCU 或 benchmark。
\end{record}

\end{document}
